% Study 13. Read by main.tex only when generated-native-n1/evidence.tex exists.
% Every result below is a macro or table written by media_rl.native_actuator_qualification_report
% from the registered final analysis; the outcome sentences were fixed there before the interim.
\section{Study 13: does the requested ratio move a native encoder?}
\label{sec:native-n1}
Every controller in Studies 10 to 12 acts by requesting a bitrate as a ratio of the bandwidth estimate, and every result there comes from a fluid simulator in which the encoder follows the request. This study tests that one link of the chain on a native browser stack and nothing else. It fits no model, lets no learned controller act, and makes no claim about delivered quality or about which ratio is better in which state. Its question, estimator, test, stopping rule and the code for all of them were registered and hash-frozen before any of its data existed (\path{docs/NATIVE_PHASE_N1_ACTUATOR_QUALIFICATION.md}, tag \texttt{native-n1-preregistered}). The sentences below that state the outcome were fixed in the evidence exporter before the interim look, one per registered outcome.

\paragraph{Why it comes first.}
A powered native panel for a rate-to-quality effect is registered separately and needs 96 independently licensed sources. If three requested ratios do not move the encoder apart, that acquisition is not worth making, so the actuator question is asked first, on sources of its own.

\paragraph{Protocol.}
The harness is that of Appendix~\ref{app:native-frames}: two peers in one headless Chrome negotiate VP8 with transport-wide congestion feedback, and all media crosses an owned userspace UDP relay with a 30{,}000-byte drop-tail queue and 25\,ms of propagation delay per direction. Each fresh peer pair plays a 20-second window of a recorded source, drawn to a $640\times360$ canvas at 30\,Hz, for 18.6\,s under one of four capacity schedules: underutilised (4\,Mbps throughout), near capacity (0.85\,Mbps), queue building ($1.2\to0.6\to0.6$\,Mbps) and collapse with recovery ($1.2\to0.35\to1.2$\,Mbps). Six times per peer the sender's bandwidth estimate is read, and for the next 3\,s the encoder's maximum bitrate is held at 0.65, 0.85 or 1.05 times that value. The three arms are drawn independently with probability one third each from a seed fixed in the sealed plan, so every assignment was determined before capture. The outcome of a hold is the encoded video payload rate before packetisation, over the window from 1.8 to 2.6\,s after the request is acknowledged, divided by the estimate the request was based on. Holds with zero output, with an aliased or clipped request, or without a plateau stay in the analysis as assigned. The models of the earlier native studies run in shadow and never act.

\paragraph{Unit of inference and sample.}
The unit is the source: files sharing a title, a capture group or media lineage form one cluster, which contributes one 20-second window and one pair of numbers. Peers and holds are repeated measurements within a source, ten fresh peers per schedule and six holds per peer, 240 assigned holds per source. Sources are Wikimedia Commons videos under CC BY 3.0, CC BY 4.0 or CC0 1.0. A seeded queue of 36 candidates was fetched and reviewed for rights against archived licence evidence; the reviewer of record is an AI system, and the repository owner adopted the review as the importer's assertion. Twenty-two candidates passed and 21 were decodable. One screening of those 21, on frame statistics alone, fixed three complexity terciles of seven sources each, and a seeded draw fixed their order: twelve first-stage sources, four per tercile, and nine in reserve.

\paragraph{Estimand and test.}
Per source, the mean normalised rate is taken per schedule and arm over complete holds, the four schedules are averaged with equal weight, and the two adjacent differences are formed: 0.85 minus 0.65 and 1.05 minus 0.85. The estimand is the mean of each difference over sources with the three terciles weighted equally. It is an intention-to-treat effect of the assigned ratio, with no model and no covariate adjustment. A requested step is 0.20 of the estimate, so a difference of 0.10 means that half the step was realised. The registered margin is \NOneMargin{}. The outcome is \emph{qualified} if the one-sided 97.5\% lower bounds of both differences exceed the margin, \emph{not separable at the registered threshold} if either one-sided 98.75\% upper bound is below it, and \emph{inconclusive} otherwise or if a capture-support gate fails. The standard error is that of Stein's two-stage procedure: the standard deviation of the per-source differences over the first \NOneFirstStage{} sources, with \NOneDegrees{} degrees of freedom, whatever the final number of sources.

\paragraph{Blinded sample size.}
After the first \NOneFirstStage{} sources a routine computes the variance across sources of the per-source differences and releases one number, the prescribed total of 12, 15, 18 or \NOneMaximum{} sources, together with outcome-free capture-support figures. The variances go to a sealed record that only the final analysis opens. No mean, difference or arm-labelled figure is released, and the capture step computes no outcome statistic. The custodian of the sealed record is the repository owner, who also ran the analysis, so the blinding is procedural and not independent.

\subsection{Result}
The capture ran on one Apple M1 laptop with \NOneBrowser{}. It collected \NOnePeers{} fresh peer pairs, and the analysis covers \NOneCohorts{} assigned holds on \NOneSources{} sources (Table~\ref{tab:n1-sources}). \NOneInterimSentence{} Capture support \NOneSupport{}: \NOneComplete\% of assigned holds are complete (at least 95\% required), \NOneAlias\% are aliased or clipped (at most 10\%), completeness differs between arms by \NOneGap{} percentage points (at most 5), the test of arm shares against one third each gives $p$ of \NOneArmP{} (at least 0.001), and \NOneSkipped{} sources were passed over as not evaluable (at most 4).

\begin{table}[t]
\centering\small
\resizebox{\linewidth}{!}{\begin{tabular}{llrrrrr}
\toprule
Endpoint & Contrast of assigned ratios & Mean & Lower bound & Upper bound & $p$ & Share of step (\%) \\
\midrule
Encoded rate / frozen BWE (confirmatory) & 0.85 $-$ 0.65 & $+$0.144 & $+$0.128 & $+$0.162 & $<$0.001 & 72 \\
 & 1.05 $-$ 0.85 & $+$0.094 & $+$0.067 & $+$0.125 & $<$0.001 & 47 \\
RTP send rate / frozen BWE (secondary) & 0.85 $-$ 0.65 & $+$0.145 & $+$0.131 & $+$0.163 & $<$0.001 & 73 \\
 & 1.05 $-$ 0.85 & $+$0.092 & $+$0.066 & $+$0.124 & $<$0.001 & 46 \\
\bottomrule
\end{tabular}
}
\caption{Native actuator qualification. Mean over \NOneSources{} sources of the difference between adjacent assigned ratios in rate divided by the frozen bandwidth estimate, terciles weighted equally. Lower bounds are one-sided 97.5\% and upper bounds one-sided 98.75\%, both with the first-stage standard error on \NOneDegrees{} degrees of freedom; $p$ tests the difference against the margin of \NOneMargin{}. The last column is the mean as a share of the requested step of 0.20. The send-rate rows are secondary and do not enter the outcome.}
\label{tab:n1-primary}
\end{table}

The difference between the two lower ratios is \NOneLowerMean{} with lower bound \NOneLowerLow{}, and between the two upper ratios \NOneUpperMean{} with lower bound \NOneUpperLow{} (Table~\ref{tab:n1-primary}): the encoder realises \NOneLowerShare\% and \NOneUpperShare\% of the requested steps. \NOneOutcomeSentence{} \NOneConsequenceSentence{}

\paragraph{Secondary analyses.}
None of these changes the outcome. The RTP send rate gives \NOneSendLowerMean{} and \NOneSendUpperMean{}, which the same rule would read as \NOneSendVerdict{}. Within each capacity schedule the same two differences are tested with Holm's correction over the \NOneRegimeTests{} tests (Table~\ref{tab:n1-regimes}); at a family-wise one-sided level of 0.025, \NOneRegimeEstablished{} of them are established above the margin, and what is not is \NOneRegimeOpen{}. By complexity tercile the intervals are descriptive and are not adjusted for the two-stage design (Table~\ref{tab:n1-complexity}). \NOneDesignationSentence{} Of the \NOneSources{} sources, \NOneSeparating{} have both of their own differences above the margin, which gives an exact one-sided 97.5\% lower bound of \NOneSeparatingBound\% on that share. The standard deviation across sources of a per-source difference is \NOneLowerSd{} \NOneLowerSdInterval{} and \NOneUpperSd{} \NOneUpperSdInterval{}; the size of any later source panel depends on this quantity. The source-clustered first-stage $F$ is \NOneEncoderF{} with partial $R^2$ \NOneEncoderRsq{}, which are instrument-strength indicators for a later quality stage. A plateau was observed in \NOnePlateauLow\%, \NOnePlateauMid\% and \NOnePlateauHigh\% of holds at the three ratios.

\begin{table}[t]
\centering\small
\begin{tabular}{llrrrr}
\toprule
Network regime & Contrast & Mean & Lower bound & Upper bound & Holm $p$ \\
\midrule
underutilized & 0.85 $-$ 0.65 & $+$0.174 & $+$0.144 & $+$0.210 & $<$0.001 \\
 & 1.05 $-$ 0.85 & $+$0.145 & $+$0.112 & $+$0.184 & $<$0.001 \\
near-capacity & 0.85 $-$ 0.65 & $+$0.152 & $+$0.126 & $+$0.182 & $<$0.001 \\
 & 1.05 $-$ 0.85 & $+$0.059 & $+$0.033 & $+$0.090 & 0.031 \\
queue-building & 0.85 $-$ 0.65 & $+$0.153 & $+$0.128 & $+$0.183 & $<$0.001 \\
 & 1.05 $-$ 0.85 & $+$0.044 & $+$0.013 & $+$0.080 & 0.169 \\
collapse-recovery & 0.85 $-$ 0.65 & $+$0.095 & $+$0.053 & $+$0.144 & 0.011 \\
 & 1.05 $-$ 0.85 & $+$0.126 & $+$0.056 & $+$0.208 & 0.017 \\
\bottomrule
\end{tabular}

\caption{Secondary: the two differences of Table~\ref{tab:n1-primary} within each capacity schedule, with the same bounds and Holm-corrected $p$ over the eight tests.}
\label{tab:n1-regimes}
\end{table}

\begin{table}[t]
\centering\small
\begin{tabular}{lrll}
\toprule
Complexity tercile & Sources & 0.85 $-$ 0.65 & 1.05 $-$ 0.85 \\
\midrule
low & 4 & $+$0.131 [$+$0.103, $+$0.160] & $+$0.073 [$+$0.027, $+$0.118] \\
medium & 4 & $+$0.146 [$+$0.118, $+$0.175] & $+$0.089 [$+$0.044, $+$0.135] \\
high & 4 & $+$0.153 [$+$0.124, $+$0.182] & $+$0.119 [$+$0.073, $+$0.164] \\
\bottomrule
\end{tabular}

\caption{Diagnostic: the two differences by complexity tercile, with 95\% intervals from the pooled within-tercile variance. The intervals are descriptive, and a tercile mean alone supports no conclusion.}
\label{tab:n1-complexity}
\end{table}

\begin{table}[t]
\centering\small
\resizebox{\linewidth}{!}{\begin{tabular}{rrlll}
\toprule
Order & Commons page & Source title & Licence & Complexity tercile \\
\midrule
1 & 185002032 & OpenRefine3-10-presentation & CC-BY-4.0 & low \\
2 & 70401610 & G.E.M. [\ldots] MV (HD) [\ldots] & CC-BY-3.0 & medium \\
3 & 194480872 & 8200 [\ldots] & CC0-1.0 & high \\
4 & 78181584 & Joy \& Heron - Animated CGI Spot by Passion Pictures & CC-BY-4.0 & low \\
5 & 146207667 & Castejón - Gamazada urteurrena (2024) (16) & CC-BY-4.0 & medium \\
6 & 42522829 & WV0016-D13 FerteSaintCyr & CC-BY-3.0 & high \\
7 & 29832113 & Paroxysm at Etna, 16-17 November 2013 & CC-BY-3.0 & low \\
8 & 98763505 & Moscow trolleybus 9803 2021-01 GX010217 & CC0-1.0 & medium \\
9 & 34935632 & [\ldots] - [\ldots] -110. [\ldots] & CC-BY-3.0 & high \\
10 & 187234232 & Entrevista a Edson Canaza - 12° Encuentro de cine Corriente. webm & CC-BY-4.0 & low \\
11 & 45099858 & Demolició Nou Ajuntament & CC-BY-4.0 & medium \\
12 & 115782037 & Brajkovič na startu kronometra na Dirki po Sloveniji 2012 & CC-BY-3.0 & high \\
\bottomrule
\end{tabular}
}
\caption{The analysed sources in the pre-drawn order, with their Wikimedia Commons page numbers. Title text in scripts the paper's font does not cover is elided as [\ldots]; full titles, authors, URLs, licence evidence and file hashes are recorded under \texttt{acquisition/native-n1/}.}
\label{tab:n1-sources}
\end{table}

\paragraph{Reading the result.}
This paragraph was written after the outcome was known and is descriptive. The qualification is an average over schedules and sources, and two features of it bear on controller design. First, the encoder follows a request only in part: it realises \NOneLowerShare\% of the lower step and \NOneUpperShare\% of the upper one, whereas the simulator of Studies 10 to 12 has the encoder take up any requested rate exactly, 1.5\,s after an increase and 0.1\,s after a decrease. Second, within single schedules what is not established above the margin after correction is \NOneRegimeOpen{}; these are the two schedules in which the link stays constrained, at 0.85\,Mbps throughout and at 0.6\,Mbps for the last two thirds of a session. Both features are consistent with the nature of the actuator. The request is a cap on the encoder's maximum bitrate, and a cap cannot raise the encoder above the rate the browser's own congestion controller allots, so a requested ratio above one is not expected to be realised in full. N1 does not separate that ceiling from encoder undershoot. The option sets of Studies 10 to 12 include a ratio above one, and their simulated sender follows it; this cap offers no such mechanism.
\ifdefined\NOneExpLevelLow
\subsection{Exploratory description of the same holds}
\label{sec:native-n1-exploratory}
This subsection is not part of the registration and supports no confirmatory claim. Its quantities were written down and committed after the final analysis had been read and before any of them was computed (\path{docs/NATIVE_PHASE_N1_EXPLORATORY.md}). The unit is the source, intervals are two-sided 95\% $t$ intervals over the \NOneSources{} sources, and they are not corrected for the number of quantities reported.

\begin{table}[t]
\centering\small
\resizebox{\linewidth}{!}{\begin{tabular}{lrrrlll}
\toprule
 & \multicolumn{3}{c}{Encoded rate / estimate at ratio} & Log growth of estimate & On-time PSNR & On-time fraction \\
Schedule & 0.65 & 0.85 & 1.05 & 1.05 $-$ 0.85 & 1.05 $-$ 0.65 & 1.05 $-$ 0.65 \\
\midrule
all, equal weight & 0.59 & 0.73 & 0.83 & $+$0.18 [$+$0.12, $+$0.24] & $+$1.3 [$+$0.3, $+$2.2] & $+$0.02 [$-$0.01, $+$0.06] \\
underutilized & 0.60 & 0.78 & 0.92 & $+$0.48 [$+$0.44, $+$0.51] & $+$2.3 [$+$1.1, $+$3.6] & $+$0.05 [$+$0.01, $+$0.10] \\
near-capacity & 0.61 & 0.76 & 0.82 & $+$0.07 [$-$0.05, $+$0.19] & $+$1.3 [$-$0.3, $+$3.0] & $+$0.03 [$-$0.02, $+$0.08] \\
queue-building & 0.58 & 0.74 & 0.78 & $-$0.00 [$-$0.09, $+$0.09] & $+$0.5 [$-$1.5, $+$2.6] & $+$0.00 [$-$0.06, $+$0.07] \\
collapse-recovery & 0.56 & 0.65 & 0.78 & $+$0.18 [$+$0.02, $+$0.34] & $+$0.8 [$-$1.4, $+$3.0] & $+$0.01 [$-$0.06, $+$0.09] \\
\bottomrule
\end{tabular}
}
\caption{Exploratory. Per capacity schedule: the encoded rate as a share of the frozen estimate at each assigned ratio; the difference between ratios 1.05 and 0.85 in the log growth of the estimate from one hold to the next; and the difference between ratios 1.05 and 0.65 in the on-time sampled PSNR contribution and in the on-time fraction of the frames captured in the outcome window. Brackets are 95\% intervals over sources.}
\label{tab:n1-exploratory}
\end{table}

\paragraph{Where the encoder sits.}
At assigned ratios of 0.65, 0.85 and 1.05 the encoded rate is \NOneExpLevelLow{}, \NOneExpLevelMid{} and \NOneExpLevelHigh{} of the frozen estimate (Table~\ref{tab:n1-exploratory}). The encoder runs somewhat below each cap, and at the top ratio it stays below the estimate itself: at \NOneExpFreeLevelHigh{} under the unconstrained schedule and between \NOneExpHeldLevelMin{} and \NOneExpHeldLevelMax{} under the three others.

\paragraph{The cap moves the estimate.}
The arm of a hold is randomised independently of everything before it, so its effect on the estimate from which the next hold starts is identified. Under the unconstrained schedule the estimate grows over one hold by \NOneExpFreeGrowthLow\%, \NOneExpFreeGrowthMid\% and \NOneExpFreeGrowthHigh\% at the three ratios, and the difference in log growth between 1.05 and 0.85 is \NOneExpFreeGrowthTop{} \NOneExpFreeGrowthTopT{}. Over all schedules that difference is \NOneExpGrowthTop{} \NOneExpGrowthTopT{}, and between 0.85 and 0.65 it is \NOneExpGrowthLow{} \NOneExpGrowthLowT{}. A cap below the estimate holds the estimate back when the link has room. On this stack a ratio of the estimate is therefore not a fixed reference: a low request now lowers the base to which later requests are relative.

\paragraph{A pilot look at delivered quality.}
Each hold also carries the on-time sampled PSNR contribution of the frames captured in its outcome window. This is the proxy that the powered panel registers, and it is not perceptual quality. Over all schedules the top ratio scores \NOneExpUtility{} \NOneExpUtilityT{} above the bottom one and \NOneExpUtilityTop{} \NOneExpUtilityTopT{} above the middle one. Under the unconstrained schedule the top-against-bottom difference is \NOneExpFreeUtility{} \NOneExpFreeUtilityT{}, and \NOneExpFreeBestTop{} of the \NOneSources{} sources score best at the top ratio. Under the \NOneExpHeldSchedules{} other schedules its interval excludes zero in \NOneExpHeldUtilityEstablished{}, and a difference between schedules is not established (the smallest Holm-adjusted $p$ of the three interaction tests is \NOneExpInteractionP{}). Of the \NOneExpUtilityContrasts{} schedule-level contrasts, \NOneExpLowerBetter{} favour the lower ratio with an interval that excludes zero. Split by pairs of holds (Table~\ref{tab:n1-exploratory-holds}), the top ratio is ahead in the first two holds of a session under \NOneExpEarlyAhead{} of the \NOneExpSchedules{} schedules; of the \NOneExpLaterCells{} later cells, \NOneExpLaterEstablished{} have an interval that excludes zero and \NOneExpLaterLowerBetter{} favour the lower ratio.

\begin{table}[t]
\centering\small
\begin{tabular}{llll}
\toprule
Schedule & Holds 1--2 & Holds 3--4 & Holds 5--6 \\
\midrule
underutilized & $+$5.1 [$+$2.3, $+$8.0] & $+$1.2 [$+$0.2, $+$2.1] & $-$0.3 [$-$1.8, $+$1.2] \\
near-capacity & $+$3.5 [$+$1.4, $+$5.6] & $-$1.0 [$-$4.6, $+$2.7] & $-$0.4 [$-$3.1, $+$2.3] \\
queue-building & $+$3.3 [$+$0.3, $+$6.4] & $-$2.2 [$-$5.4, $+$1.0] & $+$1.4 [$-$3.5, $+$6.3] \\
collapse-recovery & $+$4.4 [$+$1.5, $+$7.2] & $-$0.6 [$-$2.8, $+$1.6] & $-$0.4 [$-$4.8, $+$4.1] \\
\bottomrule
\end{tabular}

\caption{Exploratory. Difference between ratios 1.05 and 0.65 in the on-time sampled PSNR contribution, by capacity schedule and pair of holds within a session, with 95\% intervals over sources. The pairs fall mostly in the first, second and third capacity phase.}
\label{tab:n1-exploratory-holds}
\end{table}

\ifdefined\NOneCondSignals
\paragraph{A second pass, by what the sender had seen.}
A controller does not observe the schedule. A second pass over the same rows, with its design and its rule for naming a candidate again fixed before it was computed, asks whether backing off is better in the holds that a signal observed before the hold would flag: the top third of round-trip time, of its trend, of loss and of jitter, and the bottom third of the estimate's growth since the previous hold (Table~\ref{tab:n1-conditional}). The benefit of backing off is the proxy at ratio 0.65 minus the proxy at 1.05. The arm is randomised independently of the state before the hold, so within the flagged holds that benefit is a causal effect; the difference between flagged and other holds describes effect modification. \NOneCondCandidateSentence{} Of the \NOneCondSignals{} signals, \NOneCondUnusable{} flag fewer than a tenth of holds and say nothing. For the other three the benefit among flagged holds is \NOneCondGrowthFlagged{} \NOneCondGrowthFlaggedT{}, \NOneCondJitterFlagged{} \NOneCondJitterFlaggedT{} and \NOneCondRttFlagged{} \NOneCondRttFlaggedT{}, against \NOneCondGrowthRest{}, \NOneCondJitterRest{} and \NOneCondRttRest{} among the other holds, and the smallest Holm-adjusted $p$ of a difference is \NOneCondBestHolm{}. The three signals overlap, so their agreement in sign is one observation and not three. A second look at one sample that names no candidate is weak evidence. What it adds is a direction and a size for a confirmatory study on fresh sources.

\begin{table}[t]
\centering\small
\resizebox{\linewidth}{!}{\begin{tabular}{llrlllr}
\toprule
Signal before the hold & Flagged when & Flagged (\%) & Benefit, flagged & Benefit, other holds & Difference & Holm $p$ \\
\midrule
growth of the estimate since the previous hold (log) & below 0.047 & 33 & $+$0.5 [$-$0.8, $+$1.8] & $-$0.8 [$-$1.8, $+$0.3] & $+$1.3 [$-$0.1, $+$2.6] & 0.27 \\
jitter (ms) & above 12.0 & 33 & $+$1.5 [$-$0.2, $+$3.3] & $-$0.8 [$-$2.3, $+$0.8] & $+$2.3 [$-$0.2, $+$4.9] & 0.27 \\
round-trip time (ms) & above 57.0 & 33 & $+$1.0 [$-$0.7, $+$2.7] & $-$0.9 [$-$2.4, $+$0.5] & $+$2.0 [$-$0.6, $+$4.5] & 0.27 \\
round-trip time trend (ms) & above 0.000 & 5 & too few holds & $+$0.1 [$-$1.1, $+$1.3] & -- & -- \\
loss fraction & above 0.000 & 0 & too few holds & $+$0.0 [$-$1.0, $+$1.1] & -- & -- \\
\bottomrule
\end{tabular}
}
\caption{Exploratory, second pass. Benefit of backing off, defined as the on-time sampled PSNR contribution at ratio 0.65 minus that at ratio 1.05, among holds flagged by a signal observed before the hold and among the other holds, with 95\% intervals over sources and Holm-adjusted $p$ for their difference. Thresholds are pooled tercile cut-offs.}
\label{tab:n1-conditional}
\end{table}
\fi

\paragraph{What the pilot suggests.}
By schedule, no lower ratio is shown to be better, and where the ratios differ the higher one wins, most clearly while a session is starting and when the link has room. By what the sender had seen, the estimates lean the other way in the flagged third of holds and are not established. A controller that backs off below the estimate would have to earn its gain in those moments, against a browser that already adapts its own rate, and \NOneSources{} sources cannot say whether it would. This is a reason to measure that gain on fresh sources before building on it. It is not evidence that the gain is present, and not evidence that it is absent. That measurement is registered as a separate confirmatory phase with one flag, one endpoint and a fixed sample (\path{docs/NATIVE_PHASE_B1_BACKOFF_BENEFIT.md}). It has not been run, and no source for it has been acquired.
\fi

\ifdefined\TypedHeadFamilies
\subsection{What a cap leaves of the simulator's headroom}
\label{sec:typed-headroom}
The simulated sender of Studies 10 to 12 follows ratios up to 1.05 of the estimate, and the fixed ratio selected on training traces is 1.05. The native actuator cannot exceed the estimate. A model-free check in the simulator asks how much of its hindsight headroom depends on the options that a cap cannot reach (\path{docs/JEVBWE_TYPED_HEADROOM_V1.md}). Its design and its reading rule were fixed before the run. It is development evidence, because both population panels had been read, and it trains and evaluates no model. The option set is cut off at \TypedHeadCapTop{} and at \TypedHeadLowTop{}. For each set the fixed ratio is re-selected on the tune traces, and every fixed ratio and the hindsight oracle are scored on identical traces of the two panels of \TypedHeadFamilies{} families.

\begin{table}[t]
\centering\small
\resizebox{\linewidth}{!}{\begin{tabular}{rrlrrlrl}
\toprule
Top option & Fixed ratio & Panel & Fixed & Oracle & Headroom & Share (\%) & Oracle $-$ reference fixed \\
\midrule
1.05 & 1.05 & replication & 0.082 & 0.192 & $+$0.110 [$+$0.083, $+$0.137] & 100 & $+$0.110 [$+$0.083, $+$0.137] \\
 &  & sealed & 0.174 & 0.296 & $+$0.122 [$+$0.092, $+$0.153] & 100 & $+$0.122 [$+$0.092, $+$0.153] \\
1.00 & 1.00 & replication & 0.087 & 0.172 & $+$0.085 [$+$0.062, $+$0.109] & 78 & $+$0.090 [$+$0.064, $+$0.115] \\
 &  & sealed & 0.175 & 0.276 & $+$0.101 [$+$0.075, $+$0.128] & 83 & $+$0.103 [$+$0.074, $+$0.132] \\
0.90 & 0.90 & replication & 0.048 & 0.089 & $+$0.040 [$+$0.025, $+$0.055] & 37 & $+$0.007 [$-$0.024, $+$0.038] \\
 &  & sealed & 0.140 & 0.202 & $+$0.062 [$+$0.042, $+$0.082] & 51 & $+$0.028 [$-$0.003, $+$0.059] \\
\bottomrule
\end{tabular}
}
\caption{Hindsight headroom in the simulator when the option set is cut off. Headroom is the hindsight oracle minus the fixed ratio selected for that set; the share is of the reference set's headroom; the last column is the truncated oracle against the reference set's fixed ratio (for the reference set, its own headroom). Brackets are family-level $t$ intervals.}
\label{tab:followup-headroom}
\end{table}

With the top option at the estimate a fixed policy loses nothing: its utility changes by \TypedHeadCapRepFixedChange{} \TypedHeadCapRepFixedChangeT{} and \TypedHeadCapSealFixedChange{} \TypedHeadCapSealFixedChangeT{} on the two panels. The oracle loses \TypedHeadCapRepOracleChange{} and \TypedHeadCapSealOracleChange{}, and \TypedHeadCapRepShare\% and \TypedHeadCapSealShare\% of the hindsight headroom remain (Table~\ref{tab:followup-headroom}). The rule fixed in advance reads this as the opportunity lying \TypedHeadCapRepReading{} the removed option on the first panel and \TypedHeadCapSealReading{} it on the second. What hindsight offers a controller in this simulator is therefore mostly the chance to back off at the right moment, which a cap can do, and not the option above the estimate. Cut off at \TypedHeadLowTop{}, the set keeps \TypedHeadLowRepShare\% and \TypedHeadLowSealShare\% of the headroom and the fixed policy loses \TypedHeadLowRepFixedChange{} and \TypedHeadLowSealFixedChange{}. That set also holds the sender under the application's minimum rate whenever the estimate sits at its floor, so it bounds the cost of a lower ceiling and does not measure it cleanly.

This check and the native pilot point in different directions, and they are not reconciled. In the simulator the hindsight opportunity lies in backing off; in the pilot no schedule showed a lower ratio to be better. Both can hold. A hindsight oracle picks moments, whereas the pilot averaged over schedules; the simulator's estimator is a simple bounded one, and the browser's own may already take the backoff that the oracle exploits; and the simulated encoder takes up a request exactly, which the native one did not. A state-conditional native measurement would say which of these explains the difference.
\fi

\subsection{What Study 13 shows and does not show}
This is the only study here whose outcomes come from a real encoder and a real bandwidth estimator, and it is confirmatory: one sample, analysed once by rules fixed beforehand. Its scope is narrow. It concerns the actuator through which a ratio-of-estimate controller would act and not any controller. Nothing learned was the actor, no quality was measured and no state-dependent preference between ratios was tested, so it neither supports nor undermines the simulator gains of Studies 10 to 12. The network is a userspace relay on one machine, not a kernel emulator or an Internet path. The stack is one build of one browser with one codec at one resolution. The sources are \NOneSources{} Commons videos chosen by a seeded queue and a rights screen, not a representative corpus. The blinding was procedural, with one person as custodian and analyst, and the rights review rests on an AI reviewer's reading of archived evidence that the importer adopted. The effect is that of a 3\,s hold measured from 1.8 to 2.6\,s after the request and says nothing about shorter or longer holds. Whether a learned controller improves delivered quality on a real network remains untested.
